Here is an improved `section_1_proof.tex` that avoids the unjustified “lowest terms” step by choosing a representation with minimal positive denominator, and explicitly justifies the parity implication.

```latex
\begin{proof}
Recall that an integer \(m\) is even if there exists an integer \(t\) such that
\(m=2t\).

We first justify the parity fact that will be used twice: if \(n^2\) is even, then
\(n\) is even.

Let \(n\in \mathbb{Z}\). By the division algorithm with divisor \(2\), there exist
integers \(q\) and \(r\), with \(r\in\{0,1\}\), such that
\[
n=2q+r.
\]
If \(n\) is not even, then \(r\neq 0\), since \(r=0\) would give
\(n=2q\). Hence \(r=1\), so \(n=2q+1\). Then
\[
n^2=(2q+1)^2=4q^2+4q+1=2(2q^2+2q)+1.
\]
An integer of the form \(2s+1\) is not even: if \(2s+1=2t\) for some
integer \(t\), then
\[
1=2(t-s),
\]
which is impossible for an integer \(t-s\), because no integer multiplied by
\(2\) equals \(1\). Thus, if \(n\) is not even, then \(n^2\) is not even.
Therefore, if \(n^2\) is even, then \(n\) must be even; otherwise \(n\) not
being even would force \(n^2\) not to be even.

Now suppose, for the sake of contradiction, that \(\sqrt{2}\) is rational.
Then there exist integers \(p\) and \(q\), with \(q\neq 0\), such that
\[
\sqrt{2}=\frac{p}{q}.
\]
Consider the set
\[
S=\left\{ d\in \mathbb{Z}_{>0} : \text{there exists } c\in\mathbb{Z}
\text{ such that } \sqrt{2}=\frac{c}{d} \right\}.
\]
This set is nonempty: if \(q>0\), then \(q\in S\), while if \(q<0\), then
\(-q\in S\) because
\[
\frac{p}{q}=\frac{-p}{-q}.
\]
By the well-ordering principle, since \(S\) is a nonempty subset of the positive
integers, \(S\) has a least element. Let \(b\) be this least element, and choose
an integer \(a\) such that
\[
\sqrt{2}=\frac{a}{b}.
\]
Thus \(b>0\), and \(b\) is minimal among all positive integer denominators of
representations of \(\sqrt{2}\) as a ratio of integers.

Squaring both sides gives
\[
2=\frac{a^2}{b^2}.
\]
Since \(b\neq 0\), multiplying by \(b^2\) yields
\[
a^2=2b^2.
\]
Thus \(a^2\) is even. By the parity fact proved above, \(a\) is even, so there
exists an integer \(c\) such that
\[
a=2c.
\]
Substituting this into \(a^2=2b^2\), we get
\[
(2c)^2=2b^2.
\]
Hence
\[
4c^2=2b^2,
\]
and dividing by \(2\) gives
\[
2c^2=b^2.
\]
Therefore \(b^2\) is even. By the same parity fact, \(b\) is even, so there
exists an integer \(d\) such that
\[
b=2d.
\]
Since \(b>0\), we have \(d>0\). Also,
\[
\sqrt{2}=\frac{a}{b}=\frac{2c}{2d}=\frac{c}{d}.
\]
Thus \(d\in S\). But \(d=b/2\), so
\[
0<d<b,
\]
contradicting the minimality of \(b\) as the least element of \(S\).

Therefore the assumption that \(\sqrt{2}\) is rational is false. Hence
\(\sqrt{2}\) is irrational.
\end{proof}
```